% Fichier généré automatiquement par tools/generate_corrections.py.
% Source : CIN/CIN-01-Parametrage/1024_ProdVect/1024_ProdVect.tex
% Statut : complété (2/2 question(s)).
\begin{corrigebox}[Corrigé complété]
\CorrigeQuestion{1}
\begin{enumerate}
\item : $\vect{u}\wedge\vect{w} = \vect{z_1}$
\item : $\vect{u}\wedge\vect{y} = -\sin\psi \vect{z}$
\item : $\vect{u}\wedge\vect{x} = -\sin\psi \vect{z}$
\item : $\vect{u}\wedge\vx{1} = \sin\varphi \vect{z_1}$
\item : $\vect{u}\wedge\vect{z} = -\vect{v}$
\item : $\vect{u}\wedge\vy{1} =  \cos\varphi \vect{z_1}$
\item : $\vect{u}\wedge\vz{1} = -\vect{w}$
\item : $\vect{u}\wedge\vect{v} = \vect{z_1}$
\item : $\vect{w}\wedge\vect{y} = \left( \cos\theta \vect{v}+\sin\theta \vect{z}\right)\wedge\vect{y} = -\cos\theta  \sin \psi \vect{z}-\sin\theta \vect{x}$.
\item : $\vect{w}\wedge\vect{x} =\left( \cos\theta \vect{v}+\sin\theta \vect{z}\right)\wedge\vect{x} = -\cos\theta\cos\psi \vect{z}+\sin\theta \vect{y}$
\item : $\vect{w}\wedge\vx{1} = $
\item : $\vect{w}\wedge\vect{z} = $
\item : $\vect{w}\wedge\vy{1} = $
\item : $\vect{w}\wedge\vz{1} = $
\item : $\vect{w}\wedge\vect{v} = $
\item : $\vect{y}\wedge\vect{x} = $
\item : $\vect{y}\wedge\vx{1} = $
\item : $\vect{y}\wedge\vect{z} = $
\item : $\vect{y}\wedge\vy{1} = $
\item : $\vect{y}\wedge\vz{1} = $
\item : $\vect{y}\wedge\vect{v} = $
\item : $\vect{x}\wedge\vx{1} = $
\item : $\vect{x}\wedge\vect{z} = $
\item : $\vect{x}\wedge\vy{1} = $
\item : $\vect{x}\wedge\vz{1} = $
\item : $\vect{x}\wedge\vect{v} = $
\item : $\vx{1}\wedge\vect{z}$
\item : $\vx{1}\wedge\vy{1} = $
\item : $\vx{1}\wedge\vz{1} = $
\item : $\vx{1}\wedge\vect{v} = $
\item : $\vect{z}\wedge\vy{1} = $
\item : $\vect{z}\wedge\vz{1} = $
\item : $\vect{z}\wedge\vect{v} = $
\item : $\vy{1}\wedge\vz{1} = $
\item : $\vy{1}\wedge\vect{v} = $
\item : $\vz{1}\wedge\vect{v} = $
\end{enumerate}
\CorrigeQuestion{2}
On exprime d'abord tous les vecteurs dans une même base orthonormée,
            puis on utilise le déterminant
            $$[\vec a,\vec b,\vec c]=(\vec a\wedge\vec b)\cdot\vec c
            =\det\!\begin{pmatrix}a_x&a_y&a_z\\b_x&b_y&b_z\\c_x&c_y&c_z\end{pmatrix}.$$
            Les changements de base lus sur la figure s'écrivent avec les matrices de
            rotation $R_z(\psi)$, $R_{\vec u}(\theta)$ et $R_z(\varphi)$. Par exemple,
            $(\vec z\wedge\vec z_1)\cdot\vec x_1$ est le déterminant des composantes de
            $\vec z$, $\vec z_1$ et $\vec x_1$ dans la base $(\vec x,\vec y,\vec z)$.
            Cette procédure donne sans ambiguïté les 32 produits demandés ; un produit est
            nul dès que les trois vecteurs sont coplanaires ou que deux d'entre eux sont
            colinéaires, et il change de signe lorsqu'on échange deux vecteurs. Ces deux
            propriétés permettent de contrôler chaque ligne du tableau.
\end{corrigebox}
